\section{Graph maps and gluing across the cuts}
\label{sec:graph-gluing}

We first map each cut component to a locally finite graph using the
slab cover.  The two boundary copies of a cut surface need not have
the same image.  We will reconcile them by homotopies on small
collars and a new edge path for that surface.  Throughout this modification,
each vertex will retain its original cover set as an anchor for
the entire fiber.
This is a collar version of the bounded-width gluing argument in
\cite[Section~3.10.1, Lemma~(A)]{GromovFourLectures2023}.  We keep the
finite image subgraphs and the partition supports explicit to control
both whole fibers and local finiteness of the target.

\subsection{A locally finite partition and its graph}

Fix a component $P$ of $N$, with the cover $\mathcal U_P$ from
\eqref{eq:slab-cover}.  Choose a locally finite smooth partition of
unity $\{\psi_\alpha\}$ whose supports are compact and with each
support contained in some assigned $U(\alpha)\in\mathcal U_P$.
Such a partition follows by taking a locally finite relatively
compact refinement of the cover and then a subordinate partition
of unity.  Group functions assigned to the same set:
\begin{equation}
\label{eq:grouped-partition}
 \phi_U=\sum_{\alpha:U(\alpha)=U}\psi_\alpha.
\end{equation}
There are only finitely many summands in each of these sums.  Indeed,
their nonempty supports all meet the compact set $\overline U$,
whereas the original family of supports is locally finite.
Consequently $\supp\phi_U$ is a compact subset of $U$.  Moreover,
the family $\{\supp\phi_U\}$ remains locally finite: a neighborhood
meeting only finitely many of the original supports meets only
finitely many grouped supports.  Discard identically zero functions.

Give each remaining function a vertex $v$, and write $U_v$ and
$\phi_v$ for the corresponding set and function.  Include the edge
$[v,w]$ precisely when $\phi_v$ and $\phi_w$ are simultaneously
positive somewhere.  This defines a simplicial graph $K_P$.  It is
locally finite, because the compact support of $\phi_v$ meets only
finitely many supports of other functions.  Since the original
cover has multiplicity at most two, at most two of the functions
are positive at any point.  The formula
\begin{equation}
\label{eq:initial-graph-map}
 F_P(x)=\sum_v\phi_v(x)v
\end{equation}
therefore defines a map to $K_P$, interpreted in barycentric
coordinates on its vertices and edges.  Near each point only
finitely many coordinates occur, so this map is continuous for
the usual polyhedral topology, not merely for a topology on an
infinite coordinate product.

Use disjoint vertex sets for different components $P$, and put
$K_0=\coprod_P K_P$.  The maps $F_P$ give a continuous map
$F_0:N\to K_0$.  If $y$ belongs to the relative interior of a
simplex of $K_0$ and $v$ is any vertex of that simplex, then
\begin{equation}
\label{eq:initial-fiber-anchor}
 F_0(x)=y\quad\Longrightarrow\quad
 \phi_v(x)>0\quad\Longrightarrow\quad x\in U_v.
\end{equation}
In particular, the whole fiber is controlled by a single set
in \eqref{eq:cover-diameter}.

\subsection{Collar homotopies with controlled anchors}

Choose disjoint closed collars $\mathcal C_B$ of all boundary
components $B$ of $N$, with the family of collars locally finite.
Take them from sufficiently small tubular neighborhoods of the
cut surfaces in $M$.  Each collar is compact, and its normal
$g$-length may be taken at most one.  If $B^+,B^-$ are the two
copies of a cut surface $S$, the intrinsic diameter estimate
$\diam_g S\le2\pi$ from
Proposition~\ref{prop:separating-cut} gives
\begin{equation}
\label{eq:paired-collar-diameter}
 \diam_g q(\mathcal C_{B^+}\cup\mathcal C_{B^-})
 \le L,\qquad L:=2\pi+2.
\end{equation}
Indeed, project any two collar points normally to $S$, join the
projections in $S$, and add the two normal segments.

Write each collar as $B\times[0,1]$, with $B\times\{0\}$ the
actual boundary and $B\times\{1\}$ its inner end.  Let
\[
 h_B:B\longrightarrow K_0,\qquad h_B(x)=F_0(x,1).
\]
Its image lies in a finite connected subgraph $L_B$ of the
appropriate $K_P$: take the smallest subcomplex containing
$h_B(B)$.  Compactness of $B$ and local finiteness of $K_P$ imply
finiteness of this subcomplex, and connectedness of the image
implies connectedness of the subcomplex.  More importantly, its
choice ensures that
\begin{equation}
\label{eq:homotopy-anchor}
 v\in L_B\quad\Longrightarrow\quad
 (B\times\{1\})\cap\{\phi_v>0\}\ne\varnothing.
\end{equation}
For a vertex in the image this is immediate.  Any other vertex
of $L_B$ is an endpoint of an edge whose interior meets the
image, and both endpoint coordinates are positive at such a
point.  Thus \eqref{eq:homotopy-anchor} concerns the same cover
sets $U_v$ as the original partition, even if the ensuing
homotopy visits parts of $L_B$ outside $h_B(B)$.

Each $B$ is a sphere or a real projective plane by
Proposition~\ref{prop:separating-cut}, so $\pi_1(B)$ is finite.
Since the fundamental group of a connected graph is free and
torsion free, $(h_B)_*$ is trivial.  The map $h_B$ therefore
lifts to the universal covering tree of $L_B$.  Contracting
this lift to a lift of a chosen vertex $a_B\in L_B$ and
projecting back gives a homotopy
\begin{equation}
\label{eq:boundary-nullhomotopy}
 H_B:B\times[0,1]\longrightarrow L_B,
 \qquad H_B(x,0)=h_B(x),\quad H_B(x,1)=a_B.
\end{equation}
In particular, the homotopy stays inside a subgraph all of whose
vertices satisfy \eqref{eq:homotopy-anchor}.

\subsection{Adding edge paths and descending to the manifold}

For each cut surface $S$ add an edge path $A_S$ from $a_{B^+}$ to
$a_{B^-}$, with interior disjoint from $K_0$ and from every
other added path.  More explicitly, introduce two distinct new
vertices $c_S,d_S$ and the three edges
\[
 [a_{B^+},c_S],\qquad[c_S,d_S],\qquad[d_S,a_{B^-}].
\]
All new vertices are distinct for different $S$.  This produces
an ordinary simplicial graph even if the two endpoints agree
or other paths have the same endpoints.  Parametrize this path
by $\alpha_S:[0,1]\to A_S$, taking the first chosen endpoint
at zero and the second at one; when these endpoints agree,
the parametrization is injective on $[0,1)$ and only its
endpoints are identified.  Let $m_S=\alpha_S(1/2)$.

Call the enlarged graph $K$.  It is locally finite.  The only
point requiring verification is the number of added paths at
an old vertex $v$.  An endpoint attached there is some
$a_B=v$, and \eqref{eq:homotopy-anchor} implies that the
corresponding collar meets $\supp\phi_v$.  A compact set
meets only finitely many members of a locally finite family
of collars, proving the assertion.

Define $\widetilde F:N\to K$ to agree with $F_0$ outside
the collars.  On the homotopy half of a collar, adjacent
to its inner end and given by $1/2\le s\le1$, put
\[
 \widetilde F(x,s)=H_B(x,2-2s).
\]
For the remaining halves of the paired collars of $S$, put
\begin{equation}
\label{eq:arc-collar-map}
 \begin{aligned}
 \widetilde F(x,s)&=\alpha_S(\tfrac12-s)
       &&\text{on }\mathcal C_{B^+},\quad 0\le s\le\tfrac12,\\
 \widetilde F(x,s)&=\alpha_S(\tfrac12+s)
       &&\text{on }\mathcal C_{B^-},\quad 0\le s\le\tfrac12.
 \end{aligned}
\end{equation}
The formulas agree at $s=1/2$ and with $F_0$ at $s=1$.
Local finiteness of the collars proves continuity on all of
$N$.  Both boundary copies of $S$ are mapped constantly to
$m_S$.  Since $q:N\to M$ is the quotient that identifies
these copies, $\widetilde F$ descends to a continuous map
\begin{equation}
\label{eq:final-graph-map}
 f:M\longrightarrow K,\qquad f\circ q=\widetilde F.
\end{equation}

\subsection{The diameter of an entire fiber}

Fix $y\in K_0$, and let $v$ be a vertex of the unique
simplex whose relative interior contains $y$.  Every point
of $f^{-1}(y)$ lies at distance at most $L$ from the one
set $q(U_v)$.  To prove this, lift the point to $N$.  If
its value is unchanged from $F_0$, use
\eqref{eq:initial-fiber-anchor}.  Otherwise it lies in some
collar $\mathcal C_B$.  On the homotopy half, its image
lies in $L_B$; because $L_B$ is a subcomplex and contains
$y$, it contains $v$.  On the remaining half the image can
belong to $K_0$ only at the endpoint $a_B$, so again
$v\in L_B$.  By \eqref{eq:homotopy-anchor}, choose
$z\in B\times\{1\}$ with $\phi_v(z)>0$, hence
$z\in U_v$.  Equation~\eqref{eq:paired-collar-diameter}
gives the asserted distance to $q(z)$.

For any two points of this fiber, their anchor points
therefore belong to the same set $q(U_v)$.  It follows
from \eqref{eq:cover-diameter} that
\begin{equation}
\label{eq:whole-fiber-bound}
 \diam_g f^{-1}(y)\le440+2L
 =440+2(2\pi+2)<500.
\end{equation}
If instead $y\in K\setminus K_0$, it belongs to the
open part $A_S\setminus K_0$ of exactly one added path.  Its entire
preimage is contained in the images of the two collars
of $S$, so \eqref{eq:paired-collar-diameter} gives the
stronger bound $L$.  This also covers the new vertices
and the common value $m_S$.  Empty fibers require no
estimate.  Thus \eqref{eq:whole-fiber-bound} controls
whole fibers, without any connectedness assumption.

\begin{proof}[Proof of Theorem~\ref{thm:main}]
For the normalized spectral inequality with $\lambda=1$,
the preceding construction supplies the required graph
and map with constant $500$.  For general $\lambda>0$,
replace $g$ by $g'=\lambda g$.  Scalar curvature and
the Laplacian are multiplied by $\lambda^{-1}$, so the
spectral inequality for $g'$ has lower bound one.
Moreover $\dist_{g'}=\sqrt\lambda\,\dist_g$.
Applying the normalized conclusion and rescaling
\eqref{eq:whole-fiber-bound} gives
\[
 \diam_g f^{-1}(y)\le500\lambda^{-1/2}
 \qquad(y\in K),
\]
as claimed.
\end{proof}
